\chapter{Background}

\section{Purpose}

This chapter is a private learning layer for the mathematical background behind
intrinsic Quillen homology and cohomology of Jordan algebras. It is not a claim
registry and should not be treated as proof evidence.

\begin{note}
  Let \(k\) be a field of characteristic \(0\), and let
  \(\operatorname{Jord}_k\) denote a suitable category of Jordan algebras over
  \(k\), for example nonunital Jordan algebras, simplicial Jordan algebras, or
  dg Jordan algebras over the Jordan operad.

  The word \emph{intrinsic} means that the construction is carried out inside the
  category of Jordan algebras itself. In particular, one does not first pass from
  a Jordan algebra \(J\) to an associated Lie algebra, such as the Tits--Kantor--
  Koecher Lie algebra, and then take Lie algebra homology.

  In the nonunital absolute setting, the relevant abelianization functor is the
  functor of indecomposables
  \[
      Q(J) = J/J^2,
  \]
  where
  \[
      J^2 := \operatorname{span}_k\{xy : x,y\in J\}.
  \]
  The quotient \(J/J^2\) has zero multiplication, and it is the universal
  zero-multiplication Jordan algebra receiving a morphism from \(J\).

  To define Quillen homology, one replaces \(J\) by a cofibrant replacement
  \[
      \widetilde{J} \xrightarrow{\sim} J
  \]
  in the chosen homotopical category of Jordan algebras, and then applies \(Q\).
  Thus the derived indecomposables are
  \[
      \mathbb{L}Q(J) := Q(\widetilde{J})
      = \widetilde{J}/(\widetilde{J})^2.
  \]
  The intrinsic Jordan Quillen homology of \(J\) is
  \[
      H^{\operatorname{Jord}}_n(J)
      :=
      H_n\bigl(\mathbb{L}Q(J)\bigr).
  \]

  More generally, for a morphism of Jordan algebras
  \[
      A \longrightarrow B,
  \]
  one expects a relative cotangent complex
  \[
      \mathbb{L}^{\operatorname{Jord}}_{B/A}.
  \]
  If \(M\) is an appropriate \(B\)-module, for example a Beck module or Jordan
  module, then the intrinsic Jordan Quillen cohomology is defined by
  \[
      D^n_{\operatorname{Jord}}(B/A,M)
      :=
      H^n\operatorname{Hom}_B
      \bigl(
          \mathbb{L}^{\operatorname{Jord}}_{B/A},
          M
      \bigr).
  \]
  Dually, one may define homology with coefficients by
  \[
      D^{\operatorname{Jord}}_n(B/A,M)
      :=
      H_n\bigl(
          \mathbb{L}^{\operatorname{Jord}}_{B/A}
          \otimes_B M
      \bigr),
  \]
  or, more carefully, by using the appropriate derived tensor product.

  The low-degree interpretation should be analogous to the commutative and Lie
  cases:
  \[
      D^0_{\operatorname{Jord}}(B/A,M)
      \cong
      \operatorname{Der}_A^{\operatorname{Jord}}(B,M),
  \]
  while
  \[
      D^1_{\operatorname{Jord}}(B/A,M)
  \]
  should classify square-zero extensions of \(B\) by \(M\), subject to the precise
  choice of the coefficient category.
\end{note}

\section{Working Setting}

The default ground field is a field \(k\) of characteristic \(0\). The current
working category is usually nonunital Jordan algebras, with augmented unital
algebras treated through their augmentation ideals when that convention is
appropriate.

\begin{definition}
A Jordan algebra is a commutative algebra whose product satisfies the Jordan
identity. In this project the product is written as \(xy\) in mathematical
prose and sometimes as \(x * y\) in code-oriented notes.
\end{definition}

\section{Indecomposables}

The guiding abelianization candidate is
\[
  Q(J) = J/J^2,
\]
where \(J^2\) denotes the linear span of products \(xy\). The slogan is useful
only after the ambient category and the unital convention have been fixed.

\begin{question}
What is the cleanest first setting for \(Q(J)=J/J^2\): absolute nonunital
Jordan algebras, augmented unital Jordan algebras, or a relative category?
\end{question}

\begin{note}
  In the absolute nonunital category, the abelian objects are precisely the
  Jordan algebras with zero multiplication. Indeed, a vector space \(V\) with
  product
  \[
      uv = 0
      \qquad
      \text{for all }u,v\in V
  \]
  is automatically a Jordan algebra. Therefore the natural abelianization of a
  nonunital Jordan algebra \(J\) is not a commutator quotient, but rather
  \[
      Q(J)=J/J^2.
  \]
  Here \(J^2\) is the linear span of all products. The quotient has zero
  multiplication because every product in \(J/J^2\) is represented by an element
  of \(J^2\).

  Let \(A\) be a unital Jordan algebra over \(k\). An \emph{augmentation} of \(A\)
is a unital Jordan algebra morphism
\[
    \varepsilon : A \longrightarrow k,
\]
where \(k\) is regarded as a unital Jordan algebra with its usual
multiplication.

The pair \((A,\varepsilon)\) is called an \emph{augmented unital Jordan
algebra}. Its \emph{augmentation ideal} is
\[
    A_+ := \ker(\varepsilon).
\]

Since \(\varepsilon(1_A)=1\), the canonical inclusion
\[
    k \longrightarrow A,
    \qquad
    \lambda \longmapsto \lambda 1_A
\]
splits \(\varepsilon\) as a map of \(k\)-vector spaces. Hence there is a vector
space decomposition
\[
    A \cong k\cdot 1_A \oplus A_+.
\]

Moreover, \(A_+\) is an ideal. Indeed, if \(a\in A_+\) and \(b\in A\), then
\[
    \varepsilon(ab)
    =
    \varepsilon(a)\varepsilon(b)
    =
    0,
\]
so \(ab\in A_+\). In particular, if \(a,b\in A_+\), then \(ab\in A_+\).
Thus \(A_+\) is naturally a nonunital Jordan algebra.

The augmentation ideal is important because, for a unital algebra \(A\), the
naive quotient
\[
    A/A^2
\]
is usually useless: since \(1_A\in A\), one typically has
\[
    A^2=A.
\]
Therefore, in the augmented unital setting, the correct analogue of
indecomposables is
\[
    Q(A)
    :=
    A_+/A_+^2,
\]
where
\[
    A_+ = \ker(\varepsilon).
\]
This is the same idea as in the nonunital setting, but applied to the
augmentation ideal rather than to the whole unital algebra.

Conversely, if \(J\) is a nonunital Jordan algebra, one may form its
unitalization
\[
    J^{+} := k\oplus J.
\]
The product is defined by
\[
    (\lambda,x)(\mu,y)
    :=
    \bigl(
        \lambda\mu,
        \lambda y+\mu x+xy
    \bigr),
\]
where \(\lambda,\mu\in k\) and \(x,y\in J\). The element
\[
    (1,0)
\]
is the unit. There is a natural augmentation
\[
    \varepsilon : J^{+}\longrightarrow k,
    \qquad
    \varepsilon(\lambda,x)=\lambda,
\]
and its kernel is
\[
    \ker(\varepsilon)=0\oplus J\cong J.
\]
Thus nonunital Jordan algebras may often be studied equivalently through
augmented unital Jordan algebras by passing between
\[
    J
    \qquad\text{and}\qquad
    k\oplus J.
\]
\end{note}

\section{Modules, Extensions, And Differentials}

The coefficient objects should be understood through square-zero extensions
\(J \ltimes M\), Beck modules over \(J\), and Jordan derivations. The main
research question is whether the Beck-module convention exactly matches the
classical Jordan-module convention used in the literature. (I need to ask my supervisor about this).

The expected Jordan analogue of Kahler differentials is a universal object
\(\Omega_{B/A}\) representing Jordan derivations. The corresponding cotangent
complex \(L^{\mathrm{Jord}}_{B/A}\) is still provisional until the resolution
framework is fixed.
